% Com `keys`, o gabarito aparece na própria prova; com `gabarito` (ou
% `answerkey`), uma folha de respostas é acrescentada ao final.
\documentclass[11pt,a4paper,twocols,gabarito,margin=1.0cm,font={TeX Gyre Pagella}]{avaliacao}

\begin{document}
% -----------------------------------------------------
% Cabeçalho da atividade com logo
% -----------------------------------------------------
\universidade{Universidade de Brasília}
\faculdade{Faculdade UnB Planaltina}
\curso{Ciências Naturais}
\disciplina{Física 1}
\semestre{2026/1}
\assunto{1a. Avaliação}
\turno{Diurno}
\professor{Fulano de Tal}
\header

\begin{instrucoes}%
  \item Responda todas as questões.
  \item Utilize caneta azul ou preta.
  \item Não é permitido consulta a materiais ou colegas.
  \item Ao terminar, poste suas respostas em \href{https://gaba.sistema.pro.br}
  \item Utilize a tabela abaixo quando necessário.
\end{instrucoes}

\begin{center}
  \setlength{\tabcolsep}{6pt}%
  \begin{tabular}{llcl}
    \toprule
    Grandeza         & Tamanho (m)  & Massa (kg)   & Tempo (s)                                  \\
    \midrule
    Próton           & \(10^{-15}\) & \(10^{-27}\) & X                                          \\
    Átomo            & \(10^{-10}\) & X            & X                                          \\
    Vírus            & \(10^{-7}\)  & \(10^{-19}\) & X                                          \\
    Terra            & \(10^{7}\)   & \(10^{25}\)  & Rotação: \(10^{5}\); revolução: \(10^{7}\) \\
    Sol              & \(10^{9}\)   & \(10^{30}\)  & X                                          \\
    Via Láctea       & \(10^{21}\)  & \(10^{41}\)  & X                                          \\
    Universo visível & \(10^{26}\)  & \(10^{52}\)  & Idade: \(10^{18}\)                         \\
    \bottomrule
  \end{tabular}
\end{center}
\medskip

\begin{questao}{1,5}%
  Assinale verdadeiro (V) ou falso (F) em cada uma das sentenças abaixo.\par
  \begin{vf}
    \qitem{V} O vetor velocidade instantânea é tangente à trajetória em cada um dos seus pontos.
    \qitem{F} A aceleração de uma partícula que segue uma trajetória curva aponta sempre para o centro da curva.
    \qitem{V} Se a densidade de certo material está em g/cm$^3$, ao convertermos o valor para kg/m$^3$, temos que multiplicar por $10^3$.
    \qitem{V} Os vetores $\hat{\i}$, $\hat{\j}$ e $\hat{k}$ possuem módulo igual a 1 (um) e nenhuma dimensão. Indicam apenas direções e sentidos no espaço.
    \qitem{F} A componente $x$ de um vetor de módulo 4 que faz um ângulo de 60 graus com o sentido positivo do eixo $x$ é, aproximadamente, 3,46.
  \end{vf}
\end{questao}

\begin{questao}{1}%
  Julgue os itens abaixo, considerando os vetores $\vec{a}= 3\hat{\i}+3\hat{\j}$ e $\vec{b}= 2\hat{\i}+4\hat{\j}$, em verdadeiros (V) ou falsos (F).
  \begin{vf}
    \qitem{F} O módulo do vetor $\vec{c} = \vec{a} - \vec{b}$ é 1.
    \qitem{V} O produto escalar entre $\vec{a}$ e $\vec{b}$ é $18$.
    \qitem{V} O módulo do vetor $\vec{a}$ é $\sqrt{18}$.
    \qitem{F} Se $m = 1,5$ então $\vec{F} = m\vec{a}$ é um vetor de módulo igual a  $40,5$.
  \end{vf}
\end{questao}

\begin{obs}
  Responda às duas questões seguintes considerando a posição de uma partícula cujo movimento no plano $xy$ é dada por
  $$
    \vec{r}(t)= (t^4)\hat{\i}+(6 + 2t^2)\hat{\j}
  $$
  com $\vec{r}$ em metros e $t$ em  segundos.
\end{obs}

\begin{questao}{1}%
  A aceleração média da partícula, em metros por segundo ao quadrado, entre os instantes $t=2$\,s e $t=4$\,s é dada por:
  \begin{me}
    \qitem{}  $120\hat{\i}+4\hat{\j}$
    \qitem{X} $112\hat{\i}+4\hat{\j}$
    \qitem{}  $72\hat{\i}$
    \qitem{}  $128\hat{\i}-4\hat{\j}$
    \qitem{}  $72\hat{\i}-2\hat{\j}$
  \end{me}
\end{questao}

\begin{questao}{3}%
  A velocidade da partícula, em metros por segundo, no instante $t=2$\,s, é:
  \begin{me}[2]
    \qitem{}  $32\hat{\i}-4\hat{\j}$
    \qitem{}  $32\hat{\i}+\hat{\j}$
    \qitem{}  $128\hat{\i}+4\hat{\j}$
    \qitem{}  $96\hat{\i}-6\hat{\j}$
    \qitem{X} $32\hat{\i}+8\hat{\j}$
  \end{me}
\end{questao}

\begin{questao}{1}
  Uma partícula experimenta um movimento unidimensional cuja aceleração é descrita pela função $a(t)=3t^2$. A partícula parte do repouso a uma distância de três metros da origem do sistema de coordenadas. Assinale a alternativa que melhor representa a equação cinemática da posição do móvel:
  \begin{me}[1]
    \qitem{} $x(t)=12t + 3$
    \qitem{} $x(t)=12t$
    \qitem{X} $x(t)=\frac{t^4}{4} + 3$
    \qitem{} $x(t)=\frac{t^4}{4}$
    \qitem{} $x(t)=3t^2+3$
  \end{me}
\end{questao}

\begin{questao}[075]{1}
  O piloto de um avião deseja voar de leste para oeste. Um vento de \valor{80}{km/h} sopra do norte para o sul. Se a velocidade do avião em relação ao ar é igual a \valor{320}{km/h}, qual é o ângulo (positivo, em graus) segundo o qual o avião deve ser orientado relativamente ao vetor unitário que aponta no sentido sul-norte? Considere positivo o ângulo que tem sentido anti-horário.
\end{questao}

\begin{questao}[278]{1,5}
  Um soldado, a bordo de um helicóptero, solta uma caixa pesada, na esperança de atingir um alvo no solo. Se o helicóptero está voando horizontalmente a \valor{90}{m} acima do solo com velocidade de módulo \valor{65}{m/s}, a que distância horizontal (em metros) do alvo o piloto deve soltar a caixa? Use $g = \valor{9,82}{m/s^2}$.
\end{questao}

\begin{questao}{2}
  Considere que uma partícula se mova unidimensionalmente segundo a equação $x(t)=3-6t^2+ t^4$. Nessa equação, $x$ está em metros e $t$ em \textbf{minutos}. O movimento inicia-se em $t=0$. Faça o que se pede.
  \begin{subitens}
    \qitem{4} Determine, em \textbf{m/min}, o módulo da velocidade média nos primeiros dois minutos.
    \qitem{8} Determine, em \textbf{m/min}, o módulo da velocidade instantânea em $t=2$~min.
    \qitem{60} Determine, em segundos, o instante em que a aceleração da partícula é nula.
  \end{subitens}
\end{questao}
\end{document}
